\chapter{交换对称性、自旋与两个负号}

\section{壳层填充与自旋--统计问题}

壳层模型的回顾从与电子相关的两个性质开始。电子具有自旋二分之一，其两个无量纲自旋投影为
\begin{equation}
m_s=+\frac12,
\qquad
m_s=-\frac12,
\end{equation}
并且电子服从泡利不相容原理。这些事实相互关联，但它们并不是同一个陈述。

``同一状态''这一说法需要谨慎对待。没有两个电子可以占据同一个完整的单粒子态。然而，当两个电子的自旋态不同时，它们仍可以共享同一个空间轨道。在氦的独立电子近似中，两个电子都占据最低的空间轨道，一个自旋向上，另一个自旋向下。更准确地说，``第一个向上而第二个向下''与``第一个向下而第二个向上''这两种备选情形属于一个纠缠的双电子自旋态。

这仅仅是旧的壳层模型近似。它忽略了电子之间的相互作用，此处使用它只是为了进行简单的计数。对于锂，前两个电子以相反的自旋填满最低的空间轨道。第三个电子不能进入那些完整态中的任何一个，必须占据更高的轨道。

下一步是熟悉的轨道角动量图。具有
\begin{equation}
\ell=0,\qquad \ell=1,\qquad \ell=2
\end{equation}
的扇区分别称为\(S\)、\(P\)和\(D\)态。固定的\(\ell\)值有
\begin{equation}
2\ell+1
\end{equation}
个轨道态，因此
\begin{equation}
\ell=0\Rightarrow 1,
\qquad
\ell=1\Rightarrow 3,
\qquad
\ell=2\Rightarrow 5.
\end{equation}
库仑问题具有额外的简并：来自不同\(\ell\)扇区的径向能级可以具有相同的能量。对于壳层填充的论证，下一个壳层包含一个\(S\)轨道和三个\(P\)轨道。因此它有
\begin{equation}
1+3=4
\end{equation}
个空间态。自旋使计数加倍：
\begin{equation}
2(1+3)=8.
\end{equation}

这八个完整的单电子态在从锂到氖的简单递进中被填充。在闭壳层之后，钠再次具有一个外层电子，这就是为什么这种粗略图像使它类似于锂。化学并不是目的。其要点在于半整数自旋与不相容性同时出现。

光子提供了对照。它们具有整数自旋，并且不服从不相容规则。经典的电磁波对应于许多光子占据同一个模式。玻色子可以聚集在一个共同态中，而全同费米子不能不止一次地占据一个完整态。问题在于这两种配对是否可能以其他方式存在：半整数自旋而不具有不相容性，或者整数自旋而具有不相容性。在三维空间中，自旋--统计关系说不能。这一熟悉的陈述在一维和二维空间中需要加以限定，此处不尝试对场论结果给出证明。

\begin{classroomqa}
\lecturetimestamp{00:11:48}
\audiencequestion{能够占据一个状态的玻色子量子数是否有最大值？}
\lecturerresponse{没有。玻色子没有占据数限制，而且它们倾向于聚集在同一状态中。}
\end{classroomqa}

\section{多粒子波函数}

我们将通过两个负号来接近这一联系，一个来自交换，一个来自旋转。从交换开始，并暂时略去自旋。

令\(x_i\)表示粒子\(i\)的位置。下标是粒子标签，而不是笛卡儿分量；三个空间坐标被压缩为单个符号\(x_i\)。\(n\)个粒子的位置基为
\begin{equation}
\lvert x_1,x_2,\ldots,x_n\rangle.
\end{equation}
对于抽象态\(\lvert\Psi\rangle\)，坐标波函数为
\begin{equation}
\psi(x_1,x_2,\ldots,x_n)
=
\langle x_1,x_2,\ldots,x_n\mid\Psi\rangle.
\end{equation}
其模平方
\begin{equation}
\bigl\lvert\psi(x_1,x_2,\ldots,x_n)\bigr\rvert^2,
\end{equation}
是发现粒子\(1\)在\(x_1\)、粒子\(2\)在\(x_2\)等等的联合概率密度。

现在假设这些粒子是全同的。首先想象一个电子被涂成蓝色，另一个被涂成红色。那么``蓝色在\(x_1\)，红色在\(x_2\)''就可以与``蓝色在\(x_2\)，红色在\(x_1\)区分开来。''实际的电子不带有这样的颜色或个体标签。交换记账标签不能创造出新的物理排列。

整体相位不改变概率或期望值。因此，对于处于确定交换扇区中的态，坐标波函数可以满足
\begin{equation}
\psi(x_2,x_1)=\eta\,\psi(x_1,x_2),
\qquad
\lvert\eta\rvert=1.
\end{equation}
对同一交换施加两次会恢复原来的自变量，并给出
\begin{equation}
\eta^2=1,
\end{equation}
所以两个可能的交换本征值为
\begin{equation}
\eta=+1
\qquad\hbox{or}\qquad
\eta=-1.
\end{equation}

\editorialnote{双交换代数只有在将某一状态取为具有确定的交换行为之后，才能确定可能的本征值。对于固定种类的允许态必须位于同一交换扇区这一进一步限制，是一条物理规则；它并不能仅由\(\eta^2=1\)推出。同样的限定也适用于下文的排除论证：一旦选定了费米子扇区，反对称性便蕴含排除，但交换代数本身并不选择该扇区。}

对于玻色子，允许的波函数是对称的，
\begin{equation}
\psi_B(x_1,x_2)=\psi_B(x_2,x_1),
\end{equation}
而对于费米子，它们是反对称的，
\begin{equation}
\psi_F(x_1,x_2)=-\psi_F(x_2,x_1).
\end{equation}
当负号乘在整个状态上时它是不可见的，但当振幅相加时它就有意义了。交换对称性正是以这种方式获得可观测的后果。

\section{乘积、求和与排除}

首先考虑两个粒子被分配到同一个单粒子轨道\(\psi_0\)。它们的乘积为
\begin{equation}
\psi_0(x_1)\psi_0(x_2).
\end{equation}
交换粒子坐标得到
\begin{equation}
\psi_0(x_2)\psi_0(x_1)
=
\psi_0(x_1)\psi_0(x_2).
\end{equation}
波函数值是数，不是算符，因此这些因子可以对易。同态乘积是对称的。它是一种允许的玻色子形式，但它不能描述两个全同费米子。

接下来把两个粒子分配到不同的单粒子轨道：
\begin{equation}
\psi_0(x_1)\psi_1(x_2).
\end{equation}
交换得到
\begin{equation}
\psi_0(x_2)\psi_1(x_1).
\end{equation}
对于不同的函数，交换后的乘积一般既不是原来的表达式，也不是它的负值。因此，裸乘积既不具有所要求的交换对称性。这并不禁止一个粒子占据\(\psi_0\)而另一个粒子占据\(\psi_1\)。它意味着两粒子态必须把这两种分配组合起来。

对称组合为
\begin{equation}
\psi_B^{(01)}(x_1,x_2)
\propto
\psi_0(x_1)\psi_1(x_2)
+
\psi_0(x_2)\psi_1(x_1),
\end{equation}
反对称组合为
\begin{equation}
\psi_F^{(01)}(x_1,x_2)
\propto
\psi_0(x_1)\psi_1(x_2)
-
\psi_0(x_2)\psi_1(x_1).
\end{equation}
比例符号就足够了，因为交换论证不需要归一化。这两个态都各有一个粒子处于每个轨道中；它们的不同仅在于如何组合这两种不可区分的分配。

现在让两个轨道重合。玻色子组合变为
\begin{align}
\psi_B^{(00)}(x_1,x_2)
&\propto
\psi_0(x_1)\psi_0(x_2)
+
\psi_0(x_2)\psi_0(x_1)\\
&=
2\psi_0(x_1)\psi_0(x_2),
\end{align}
它是非零的。费米子组合给出
\begin{align}
\psi_F^{(00)}(x_1,x_2)
&\propto
\psi_0(x_1)\psi_0(x_2)
-
\psi_0(x_2)\psi_0(x_1)\\
&=0.
\end{align}
在反对称交换规则下，排除随之而来：不存在两个费米子处于同一单粒子态的非零反对称态。

对于多粒子态，该规则适用于任意两个宗量的对换。一次对换使玻色子波函数保持不变，而使费米子波函数反号。重复同一次对换在两种情况下都是平凡的。令人惊讶的是，单次交换虽然不改变任何可观测的粒子标签，却改变费米子态的符号。

\section{完整粒子标签中的自旋}

现在可以通过一个小改动把自旋恢复进来。单粒子振幅同时依赖于位置和一个离散的自旋标签，
\begin{equation}
\psi(x,\sigma),
\end{equation}
其中\(\sigma=+1\)或\(-1\)标记所选的\(\sigma_z\)结果。其模平方是所述位置和自旋标签的联合概率密度。

对于多粒子，波函数为
\begin{equation}
\psi(x_1,\sigma_1;\,x_2,\sigma_2;\,\ldots).
\end{equation}
交换会交换属于两个粒子的一切，而不仅仅是它们的位置：
\begin{align}
\psi_B(x_1,\sigma_1;\,x_2,\sigma_2)
&=
+\psi_B(x_2,\sigma_2;\,x_1,\sigma_1),\\
\psi_F(x_1,\sigma_1;\,x_2,\sigma_2)
&=
-\psi_F(x_2,\sigma_2;\,x_1,\sigma_1).
\end{align}
这就是开头壳模型陈述的精确形式。两个电子只有在它们的完整状态（包括自旋）不同时，才可以共享同一个空间轨道。

\begin{classroomqa}
\lecturetimestamp{00:34:33}
\audiencequestion{交换对称性是置换群的一个实例吗？}
\lecturerresponse{是的。交换粒子是一种置换，而正号或负号的变换规律给出了置换群的一个表示。}
\end{classroomqa}

\begin{classroomqa}
\lecturetimestamp{00:35:12}
\audiencequestion{为什么两粒子构造是单粒子波函数的乘积而不是求和？}
\lecturerresponse{求和描述的是一个粒子处于叠加态。乘积属于两粒子张量积空间，描述的是每个单粒子态各分配到一个粒子。}
\end{classroomqa}

完整标签的交换如下方黑板所示。\footnote{位置--自旋标签的排序是根据可辨认的黑板画面以及来自00：33：43到00：34：25的周围讲解重建的。}
\begin{figure}[htbp]
\centering
\includegraphics[width=0.9\linewidth]{lecture_05_frame_03.png}
\caption{交换两个全同粒子会交换它们的完整位置--自旋标签。玻色子态保持不变，而费米子态改变符号。}
\lectureframe{00:35:25}
\end{figure}

这些是线性性的不同用法。对于一个粒子，
\begin{equation}
\alpha\psi_0(x)+\beta\psi_1(x)
\end{equation}
是两种备选情况的叠加。对于两个粒子，
\begin{equation}
\psi_0(x_1)\psi_1(x_2)
\end{equation}
是带有两个粒子坐标的空间上的乘积振幅。对称化或反对称化把该乘积与其交换后的配对项组合起来。

\begin{classroomqa}
\lecturetimestamp{00:36:13}
\audiencequestion{交换粒子会产生干涉效应吗？}
\lecturerresponse{会。同态费米子振幅的消失可以理解为两个交换贡献之间的相消干涉。}
\end{classroomqa}

\begin{classroomqa}
\lecturetimestamp{00:37:13}
\audiencequestion{泡利是用同样的对称性论证推导出不相容原理的吗？}
\lecturerresponse{泡利从经验事实出发，找到了一个可以表达该事实的数学框架。他最初并不是从更基本的理论推导出不相容原理的。}
\end{classroomqa}

反对称性提供了表达经验性不相容规则的数学框架。

\begin{classroomqa}
\lecturetimestamp{00:39:01}
\audiencequestion{粒子交换干涉与备选路径之间的干涉有关联吗？}
\lecturerresponse{有。一种历史是两个粒子经过而不交叉，另一种历史是它们交叉进入交换后的目的地。对于玻色子，振幅相加；对于费米子，振幅可以相消。}
\end{classroomqa}

交换提供第一个负号。旋转提供第二个。

\section{第二个负号：旋转 \texorpdfstring{\(2\pi\)}{2 pi}}

完整旋转看起来什么也没做。每个普通几何方向在转过角度 \(2\pi\) 后都回到起点。然而必须回到原状的，是概率和期望值。态矢量本身可以带一个整体相位返回，正如它在交换下可以带一个整体相位返回一样。

\begin{classroomqa}
\lecturetimestamp{00:42:12}
\audiencequestion{旋转使用哪条轴，生成元是轨道角动量、自旋，还是两者？}
\lecturerresponse{可以选择一条外部轴；下文取为 \(z\) 轴。生成元是总角动量，包括轨道和自旋两部分的贡献。自旋本身不必指向所选轴。}
\end{classroomqa}

\editorialnote{本讨论中的角动量公式使用 \(\hbar=1\)，因此 \(m\) 是 \(J_z\) 的无量纲本征值。}

令 \(\psi(\theta)\) 表示当旋转参数推进到角度 \(\theta\) 时得到的态。在本文采用的约定下，
\begin{equation}
J_z\psi(\theta)
=
-i\frac{\partial}{\partial\theta}\psi(\theta).
\end{equation}
对于具有确定 \(J_z\) 的态，
\begin{equation}
J_z\psi(\theta)=m\psi(\theta).
\end{equation}
将方程合并得到
\begin{equation}
-i\frac{\partial}{\partial\theta}\psi(\theta)
=
m\psi(\theta),
\end{equation}
其解为
\begin{equation}
\psi(\theta)=e^{im\theta}\psi(0).
\end{equation}

\(m\) 的允许值是整数或半整数。因此
\begin{equation}
\psi(2\pi)
=
\begin{cases}
+\psi(0), & m\in\mathbb Z,\\[4pt]
-\psi(0), & m\in\mathbb Z+\tfrac12.
\end{cases}
\end{equation}
再转一整圈后，
\begin{equation}
\psi(4\pi)=\psi(0)
\end{equation}
无论哪种情况都成立。这就是第二个负号：处于半整数角动量多重态中的态在 \(2\pi\) 旋转下改变符号。

该符号不改变 \(\lvert\psi\rvert^2\)，因此当它乘在整个态上时不可见。即便如此，\(2\pi\) 旋转在半整数自旋态上并不由恒等变换表示。当该符号相对于另一个相干振幅而被确定时，这一区别就变得可观测。

\section{\texorpdfstring{\(2\pi\) 与 \(4\pi\)}{2 pi 和 4 pi} 路径}

\(2\pi\) 与 \(4\pi\) 之间的区别无需更多代数就能显现。用一只手拿住杯子，不松手地将其转过一整圈。杯子回到原来的取向，但手臂被扭住了。继续转第二圈，手臂就能回到原来的位置，而无需反转杯子的运动。

一个相关的构造是把一个球放进盒子里，并用绳子把它与盒子连接起来。将球转过 \(2\pi\) 会留下一个缠结。在普通三维空间中，与 \(4\pi\) 旋转相关的缠结可以解开。相关的对象不只是球的最终取向，而是它在旋转空间中所经过的路径。

\begin{classroomqa}
\lecturetimestamp{00:54:22}
\audiencequestion{如果球绕同一轴转两次，绳子真的能解开吗？那需要另一个空间维度吗？}
\lecturerresponse{\(4\pi\) 路径可以在普通三维空间中解开。时间参数化运动；它不是另一个空间方向。\(2\pi\) 路径不能以这种方式解开，而 \(4\pi\) 路径可以。}
\end{classroomqa}

这些演示区分了非平凡的 \(2\pi\) 路径与可解开的 \(4\pi\) 路径，而无需发展旋转群的形式理论。

\section{皮带图像中的交换与旋转}

两个负号现在并排而立。交换两个全同费米子会改变它们联合波函数的符号。将一个半整数自旋态转过 \(2\pi\) 也会改变其符号。皮带给出了这两种操作之间的几何关系。

调整带子，使其两个下端已经互换。换一种看法，同一条带子带有一个\(2\pi\)扭转。将交换与\(2\pi\)旋转一起执行，会消除扭转，留下一条未扭转的带子。这一组合过程可以连续地约化为什么都不做。\footnote{交换--旋转这一观察在讲座中归功于David Finkelstein。源记录中没有指明具体出版物，因此这里仅作为该归功保留此名。}

\begin{classroomqa}
\lecturetimestamp{01:01:12}
\audiencequestion{这只是关于带子如何嵌入空间的陈述，而不是关于粒子拓扑的陈述吗？}
\lecturerresponse{这条带子可以读作时空中的粒子历史。时间向上流逝时，它的一个横截面给出类粒子对象。成对产生，一个对象旋转\(2\pi\)，并且两个相同对象被交换。在该历史中，旋转相位与交换相位相消。}
\end{classroomqa}

该历史使相位记账变得明确。从没有粒子开始。在两个分离的位置，产生一个粒子--反粒子对，在中间时刻留下四个粒子。将一个粒子旋转\(2\pi\)，同时在位形中别处交换两个相同粒子。完整历史可以形变为基本历史，其中一对粒子产生并在后来湮灭。在所讨论的符号记账中，该平凡历史没有额外相位。若\(\eta_{\mathrm{rot}}\)和\(\eta_{\mathrm{ex}}\)表示这两个符号，则所说的相消可以紧凑地写为\footnote{这个相位方程记录了从01：06：50到01：07：18所描述的相消；讲座将该构造呈现为提示性的，而不是严格证明。}
\begin{equation}
\eta_{\mathrm{rot}}\eta_{\mathrm{ex}}=1.
\end{equation}

\begin{classroomqa}
\lecturetimestamp{01:03:34}
\audiencequestion{能否通过让两个粒子各自移动经过\(\pi\)，然后再重新定向它们，把交换与旋转联系起来？}
\lecturerresponse{该构造可能反映了相关的数学，但它使用了两个分开的\(\pi\)旋转。这里所分析的过程包含一个\(2\pi\)旋转，因此这两个论证并不相同。}
\end{classroomqa}

该论证依赖于产生和湮灭类粒子对象的可能性。在一个粒子永恒存在的理论中，它不会有这种形式。它也不是作为自旋--统计定理的标准证明而提出的。其目的是展示交换符号与旋转符号彼此吻合的一个拓扑理由。

\section{团块、镜像与湮灭}

该构造可以用局域场位形重复。假设一个或多个场几乎处处取真空值，但在一个小区域中具有非零的形状和结构。该局域位形就是一个团块。

现在构造第二个位形，作为第一个的适当镜像。称一个为粒子，另一个为反粒子。将它们带到一起。场位形的对应部分相遇并相互抵消。在中间阶段，团块的某些部分可能已经消失，而其他部分仍然存在。继续形变会移除剩余结构，并使场回到真空位形。带子中两个相反排列的扭转给出同样的图景：把扭转推到一起，它们就会解开。

该过程可以反向运行。空空间可以产生一个团块及其镜像伙伴。人们可以产生两个这样的对，将一个团块旋转\(2\pi\)，交换两个相同团块，并再次湮灭这些对。同样的交换--旋转相消随后适用。这仍然是该构造的启发式推广，而不是关于每一个局域激发的定理。

\begin{classroomqa}
\lecturetimestamp{01:11:30}
\audiencequestion{这些类团块场位形有物理实现吗？}
\lecturerresponse{有。例子包括Skyrmion型描述中的核、磁单极子，以及场的其他类粒子激发。}
\end{classroomqa}

镜像团块图景使这一推广变得可理解。相关特征不是该对象最初是基本点粒子，而是一个类粒子激发及其伙伴可以被连续地产生、带到一起并湮灭。该特征促成了旋转行为与交换行为之间更广泛的关系。

\begin{classroomqa}
\lecturetimestamp{01:16:12}
\audiencequestion{自旋为二分之一是否意味着态只累积通常旋转相位的一半？}
\lecturerresponse{在此约定下相位为\(e^{im\theta}\)。对于半整数\(m\)，\(2\pi\)旋转给出\(-1\)，而\(4\pi\)旋转给出\(+1\)。}
\end{classroomqa}

\section{整体符号与零结果实验}

由\(2\pi\)旋转产生的负号能被观测到吗？先从一个观测不到它的实验开始。将电子囚禁在磁场中，让它的自旋跟随磁场，同时将装置缓慢转过\(2\pi\)，然后将电子释放到干涉装置中。将此与不做旋转的同一过程进行比较。

该旋转使态变为
\begin{equation}
\psi\longmapsto-\psi,
\end{equation}
但负号乘在整个态上。所有概率均不变，因此最终的干涉图样不保留是否发生了预备旋转的任何记录。这是一个零结果实验：一个共同相位没有任何可供干涉的对象。

\begin{classroomqa}
\lecturetimestamp{01:20:01}
\audiencequestion{为什么磁场旋转必须缓慢进行？}
\lecturerresponse{自旋必须持续跟随旋转的磁场。如果磁场拖曳得太快，自旋将无法保持锁定于其上。}
\end{classroomqa}

\section{一个电子处于两个相干分支中}

当符号被变为相对符号时，它就变得可观测了。让一个电子通过分束器。其波函数的一部分沿一条路径传播，另一部分沿另一条路径传播：
\begin{equation}
\psi(x)=\psi_1(x)+\psi_2(x).
\end{equation}
这是一个电子，而不是两个电子。它是两种备选可能性的相干叠加。必须清楚区分：这两个波包不是两个独立粒子。

将这两个波包引导到包含磁场的不同区域。在一个区域中将磁场旋转\(2\pi\)，另一个区域保持不变。对于半整数自旋态，
\begin{equation}
\psi_1(x)+\psi_2(x)
\longmapsto
\psi_1(x)-\psi_2(x),
\end{equation}
其中被旋转分支的标记被选为\(2\)。此时符号是相对的。当两个波包随后重叠时，加法已变为减法。相长干涉的一点可以变成相消干涉的一点，甚至成为节点。

\begin{classroomqa}
\lecturetimestamp{01:23:46}
\audiencequestion{如果只有一个电子却有两个波函数分支，那么被旋转的是什么？}
\lecturerresponse{电子处于两个局域化备选可能性的相干叠加中。一个分支被旋转，另一个没有。如果打开盒子以得知发生了哪一种备选可能性，该测量就会破坏干涉。}
\end{classroomqa}

打开盒子只会揭示一个电子，在一个盒子或另一个盒子中。若不做该测量，两个振幅都保留在态中。重复实验会一次一个探测地累积出干涉图样。在由相对负号产生的节点处，检测不到电子。

\section{磁进动与双缝干涉}

一种更实用的示意性安排取代了极慢的盒子实验。将相干电子波分成两条路径。在一条路径中，让电子通过一个磁场区域，其自旋初始时横向于磁场。随后自旋发生进动。通过选择磁场强度和渡越时间，该分支中的自旋可以完成整整\(2\pi\)圈，而另一分支不经历这样的旋转。\footnote{讲座指出，一个相关的旋量相位实验已经被完成，并描述了其原理，但来源记录没有标明论文或特定装置。因此这里的叙述仅限于讲座中给出的构造。}

\begin{classroomqa}
\lecturetimestamp{01:27:05}
\audiencequestion{如何确定单个电子的路径？}
\lecturerresponse{若不破坏实验，就无法确定。任何记录电子走了哪条路径的测量都会破坏干涉。}
\end{classroomqa}

\begin{classroomqa}
\lecturetimestamp{01:27:50}
\audiencequestion{波函数的分支究竟是什么，分束器又如何产生两个分支？}
\lecturerresponse{分支是沿一条路径局域化的波包。分束器将入射波变成两个空间上分离的波包，而电场或磁场可以引导这些波包沿不同路径传播。}
\end{classroomqa}

除磁场区域外，这就是双缝实验。入射电子波分成两个波包。出射波扩散到共同探测区域，在那里它们的振幅发生重叠。没有相对旋转时，中心振幅可以相加。一条路径中的\(2\pi\)自旋旋转会使一个振幅反号，并可使同样的贡献相互抵消。

\begin{classroomqa}
\lecturetimestamp{01:30:46}
\audiencequestion{如果电子穿过不同的管或狭缝，怎么会有干涉？}
\lecturerresponse{最终的探测并不揭示路径。两个出射波振幅在探测器处扩展并重叠。试图查看走的是哪条路径会破坏干涉。}
\end{classroomqa}

\begin{classroomqa}
\lecturetimestamp{01:33:22}
\audiencequestion{管或狭缝的直径必须与电子波长相当吗？}
\lecturerresponse{对于字面意义上基于衍射的双缝实现，开口必须足够小以产生所需的扩展。所描述的装置是示意性的，而非实际设计。}
\end{classroomqa}

\begin{classroomqa}
\lecturetimestamp{01:33:44}
\audiencequestion{相应的\(2\pi\)光子偏振旋转会使干涉图样保持不变吗？}
\lecturerresponse{是的。光子具有整数自旋，因此完整的\(2\pi\)旋转不会引入费米子的负号。}
\end{classroomqa}

\begin{classroomqa}
\lecturetimestamp{01:34:12}
\audiencequestion{实际实验与理想化版本之间的轻微偏离是什么？}
\lecturerresponse{理想构造使自旋绝热旋转。在实际描述中，电子穿过一个磁性区域，其自旋进动得更快，但所测量的相位仍然是\(2\pi\)圈所对应的相位。}
\end{classroomqa}

缓慢构造通过确保自旋跟随场而不产生额外运动来孤立出这一思想。进动构造以动力学方式达到相同的相对符号。在任一版本中，两条路径必须在干涉图样出现之前变得不可区分。

\begin{classroomqa}
\lecturetimestamp{01:36:15}
\audiencequestion{两条分支必须在分束器之后重新合并吗？}
\lecturerresponse{是的。出射的波包必须重叠，或扩展到共同的探测区域，以便它们的振幅能够形成干涉图样。}
\end{classroomqa}

\begin{classroomqa}
\lecturetimestamp{01:38:12}
\audiencequestion{这个单电子旋转实验是泡利不相容的结果吗？}
\lecturerresponse{不是。只有一个电子。该实验测量的是旋转下的旋量相位，而不是两个费米子不能占据同一态。}
\end{classroomqa}

\section{交换干涉与相对旋转}

交换符号可以通过同样的一般方法暴露出来，尽管所提出的装置更为复杂。从两个电子开始，将它们的联合波函数分成两条相干分支。仍然只有两个电子，而不是四个：这对电子走其中一条备选路径或另一条。在一条分支中交换两个电子，在另一条中不交换。当分支重新合到一起后，费米子交换符号表现为相对相位，并能改变干涉。

最后的提议直接比较这两个负号。一条分支获得费米子交换相位，而另一条获得相应的\(2\pi\)旋转相位。两个相位都是\(-1\)，因此相对效应抵消。\footnote{在01:40:26与01:41:15之间的记录清楚地给出了交换相位、\(2\pi\)旋转相位及其抵消，但旋转的详细归属无法从可用的讲话中恢复。} 讲座将此作为提出的思想实验来呈现，并说据演讲者所知，它尚未被实施。

\begin{classroomqa}
\lecturetimestamp{01:41:34}
\audiencequestion{观察到交换相位会挑战两个电子的不可区分性吗？}
\lecturerresponse{不会。量子不可区分性比经典意义上的相同更为微妙。它允许态矢量在交换下改变符号，即使没有任何可观测标签区分这些粒子。}
\end{classroomqa}

还有最后一点限定。将整个孤立世界旋转\(2\pi\)只会给出一个共同相位，任何内部实验都无法揭示它。可探测的实验使一条相干分支相对于另一条旋转。因此所观察到的不是绝对符号，而是由相对符号产生的干涉。
